\section{总结与延伸}

\subsection{总结表}

\begin{table}[H]
\centering
\begin{tabular}{lp{8cm}p{4cm}}
\toprule
章节 & 关注点 & 后续建议 \\
\midrule
MDP 基础回顾 & Bellman 迭代与策略/value 迭代的关系 & 复核收缩定理与演化形式 \\
表格型 Q-Learning & Off-policy 收敛与探索均衡 & 为现有任务梳理 replay buffer 的覆盖率 \\
Fitted Q-Iteration 与 DQN & 函数逼近、稳定技巧与优化细节 & 在老系统中验证 target network / double DQN 影响 \\
从 DQN 到 Actor-Critic & Continuous control 与 actor-critic 架构过渡 & 比较 deterministic policy gradient 与 Q-learning 的梯度估计 \\
系统实践与监控 & 指标、 replay buffer 健康度、干预机制 & 建立 instrumentation revoke 机制 \\
变体与成本 & Rainbow、成本模型、推理栈 & 规划训练/推理预算，并保留挂回路 \\
应用案例与治理 & benchmark 与 sim-to-real 突破、治理关注点 & 规划 soft/hard interrupt 与 anomaly 回滚 \\
\bottomrule
\end{tabular}
\caption{本讲重点内容与工程切入点}
\end{table}

\subsection{进一步阅读}

\begin{knowledgebox}{延伸资料}
\begin{itemize}
    \item Watkins \& Dayan, ``Q-Learning,'' Machine Learning 1992
    \item Mnih et al., ``Human-level Control through Deep Reinforcement Learning (DQN),'' Nature 2015
    \item Van Hasselt et al., ``Deep Reinforcement Learning with Double Q-Learning,'' AAAI 2016
    \item Lillicrap et al., ``Continuous Control with Deep Reinforcement Learning (DDPG),'' ICLR 2016
    \item Sutton \& Barto, ``Reinforcement Learning: An Introduction,'' 2nd Edition, 2018
    \item OpenAI Safety Gym 2.0 文档（可做 interruptibility 测试）
\end{itemize}
\end{knowledgebox}

\subsection{本章小结}

本讲以 MDP 与 Bellman 方程为基础，层层展开 Q-Learning 从表格到深度网络、再到 Actor-Critic 的演进，同时穿插系统实践、变体 COST 以及治理/应用案例，提供了完整的复习与操作 checklist。

\end{document}
